% !PRINTED XIII
\weberpage{XIII}

\inhaltabschnitt{Einundzwanzigster Abschnitt.}{Kreistheilungskörper.}
\inhaltpar{169}{Zerlegung der Primzahl $q$ im Kreistheilungskörper $\Omega_{q^\varkappa}$}{617}
\inhaltpar{170}{Minimalbasis des Körpers $\Omega_m$}{620}
\inhaltpar{171}{Die Primideale im Körper $\Omega_m$}{622}
\inhaltpar{172}{Die conjugirten Primideale}{623}
\inhaltpar{173}{Darstellung der Primfactoren von $p$}{628}
\inhaltpar{174}{Das \gesperrt{Kummer}'sche Theorem}{632}
\inhaltpar{175}{Die Einheitswurzeln im Körper $\Omega_m$}{638}
\inhaltpar{176}{Der in $\Omega_m$ enthaltene reelle Körper $H_m$}{641}
\inhaltpar{177}{Die Primideale des Körpers $H_m$}{643}
\inhaltpar{178}{Die Einheiten des Körpers $H_m$}{645}

\inhaltabschnitt{Zweiundzwanzigster Abschnitt.}{Abel'sche Körper und Kreistheilungskörper.}
\inhaltpar{179}{Einfache \gesperrt{Abel}'sche Körper}{648}
\inhaltpar{180}{Die Resolventen}{651}
\inhaltpar{181}{Vorbereitung zum Beweis}{654}
\inhaltpar{182}{Beweis des ersten Hülfssatzes für ein ungerades $m$}{659}
\inhaltpar{183}{Beweis des zweiten Hülfssatzes für eiu ungerades $m$}{666}
\inhaltpar{184}{Vorläufiges über den Fall eines geraden $m$}{667}

\inhaltabschnitt{Dreiundzwanzigster Abschnitt.}{Classenzahl.}
\inhaltpar{185}{Der \gesperrt{Dirichlet}'sche Satz über die Einheiten}{670}
\inhaltpar{186}{Systeme unabhängiger Einheiten und Exponentensysteme der Einheiten}{677}
\inhaltpar{187}{Fundamentalsysteme von Einheiten}{681}
\inhaltpar{188}{Reducirte Zahlen}{684}
\inhaltpar{189}{Grenzen der Anzahl der durch ein Ideal theilbaren ganzen Zahlen des Körpers $\Omega$}{686}
\inhaltpar{190}{Bestimmung des Volumens}{690}
\inhaltpar{191}{Sätze aus der Reihenlehre}{694}
\inhaltpar{192}{Anwendung auf die Bestimmung der Classenzahl}{701}

\inhaltabschnitt{Vierundzwanzigster Abschnitt.}{Classenzahl der Kreistheilungskörper.}
\inhaltpar{193}{Classenzahldarstellung im Kreistheilungskörper $\Omega_m$}{705}
\inhaltpar{194}{Bestimmung der Summen $X$}{708}
\inhaltpar{195}{Ueber die Classenzahl in dem in $\Omega_m$ enthaltenen reellen Körper}{711}
\inhaltpar{196}{Classenzahl im Körper der achten Einheitswurzeln}{717}
\inhaltpar{197}{Recurrente Berechnung der Classenzahl im Körpar $\Omega_m$, wenn $m$ eine Potenz von 2 ist}{719}
\inhaltpar{198}{Der Classenzahlfactor $A$}{722}
\inhaltpar{199}{Der Classenzahlfactor $B$}{726}
